\documentclass[11pt]{article}

%%%%%%%%%%%%%%Include Packages%%%%%%%%%%%%%%%%%%%%%%%%%%
\usepackage{xcolor}
\usepackage{mathtools}
\usepackage[a4paper, total={6in, 8in}, margin=1.25in]{geometry}
\usepackage{amsmath}
\usepackage{amssymb}
\usepackage{paralist}
\usepackage{rsfso}
\usepackage{amsthm}
\usepackage{wasysym}
\usepackage[inline]{enumitem}   
\usepackage{hyperref}
\usepackage{tocloft}
\usepackage{titlesec}
\usepackage{colortbl}
\usepackage{wrapfig}
\usepackage{pdfpages}
\usepackage[american]{circuitikz}
%%%%%%%%%%%%%%%%%%%%%%%%%%%%%%%%%%%%%%%%%%%%%%%%%%%%%%%%



\begin{document}
\Large \noindent \textbf{Physics 5C Discussion (Week 10)}\\
\normalsize
Department of Physics and Astronomy, UCLA\\
%\noindent TA: Jinyan Miao (jnmiu@ucla.edu)\\
%Office hours: Thursday 10-11 AM $\&$ 1-2 PM at PAB 2-429\\
%Tutoring center hours: Friday 2-3 PM at PAB 2-429 (Zoom room also available)\\
\noindent\rule{8cm}{0.4pt}\\
Jinyan has been addicted to playing the game called Factorio with his friends, where they build automated factories. One late-game resource is Uranium: Uranium processing mostly produces the isotope $_{92}^{238}$U, which is essential in the production of nuclear fuel and ``atomic bomb."\\

\noindent \textbf{Part (a)} In an atom of $_{92}^{238}$U, how many electrons, protons, and neutrons does it have?\\

If all of the protons, electrons, and neutrons in the $_{92}^{238}$U atom were separated, their total mass is just the sum of the individual masses, but when those particles are bind together to form the $_{92}^{238}$U atom, the atom has a different mass.\\

\noindent \textbf{Part (b)} Find the total mass of the separated protons and neutrons, from the $_{92}^{238}$U atom.\\

\noindent \textbf{Part (c)} Find the total mass of the separated electrons, from the $_{92}^{238}$U atom.\\

\noindent \textbf{Part (d)} Is the mass of the electrons important?\\

\noindent \textbf{Part (e)} Calculate the mass defect of $_{92}^{238}$U in atomic mass unit.\\

\newpage
The $_{92}^{238}$U atom is unstable, it will undergo a atomic decay chain, releasing energy in this process (and this is why we say that $_{92}^{238}$U is radioactive), until it ends up at the stable isotope $^{206}_{82}$Pb. The decay chain is a combination of $\alpha$, $\beta$, and $\gamma$ decays. \\

\noindent \textbf{Part (f)} Here are some decays in the decay chain. \\
Please help Jinyan to fill out the blank and write down the complete decay process.
\begin{enumerate}
\item The $^{238}_{92}\text{U}$ atom undergoes an $\alpha$ decay to $^{234}_{90}$Th. The process is written as $${}^{238}_{92}\text{U} \to {}^{234}_{90}\text{Th }+\alpha\,.$$
\item The $^{234}_{90}\text{Th}$ atom undergoes a $\beta$-minus decay to \underline{\qquad\quad}. This is $${}^{238}_{92}\text{U} \to \underline{\qquad\quad} +\underline{\qquad} + \bar{\nu}_e\,.$$
\item The \underline{\qquad\quad} atom undergoes a $\beta$-minus decay to $^{234}_{92}$U. This is $$
\underline{\qquad\quad} \to\ ^{234}_{92}\text{U} +\underline{\qquad} + \bar{\nu}_e\,.$$
\item The $^{234}_{92}$U atom undergoes an $\alpha$ decay to \underline{\qquad\quad}. This is
$${}^{234}_{92}\text{U} \to\underline{\qquad\quad}+\underline{\qquad}\,.$$
$$\vdots$$
\item[end.] The $^{210}_{84}$Po atom undergoes an \underline{\qquad} decay to the stable ${}^{206}_{82}$Pb. This is
$$
{}^{210}_{84}\text{Pb}\to
{}^{206}_{82}\text{U} 
+\underline{\qquad}\,.$$
\end{enumerate}


\noindent\textbf{Part (g)} One of the processes in the nuclear decay chain is a $\gamma$ decay, 
\begin{align*}
^{214}_{83}\text{Bi}^* \to {}^{214}_{83}\text{Bi} + \gamma\,.
\end{align*}
The superscript $*$ is a notation that indicates $^{214}_{83}\text{Bi}^*$ has electrons in their excited states, and the electrons "jump" to the ground state (with the corresponding atom denoted as ${}^{214}_{83}\text{Bi}$) via emitting a $\gamma$-ray. The wavelength of such $\gamma$-ray is measured to be $\lambda = 7.05\cdot 10^{-31}\,\text{m}$. Find the energy difference between the $^{214}_{83}\text{Bi}^*$ atom and the ${}^{214}_{83}\text{Bi}$ atom in this process. \\

\noindent\textbf{Part (h)} Do a quick Google search, and explain why the $\gamma$-ray released in part (g) is highly dangerous to humans. (Hint: What is the frequency of the emitted $\gamma$-ray?) Since the $\gamma$ decay process in part (g) is part of the decay chain of $^{238}_{92}$U, your answer is the main reasoning that the decay of radioactive substances like $^{238}_{92}$U can ``kill" people. \\

\noindent\textbf{Part (i)} The half-life of $^{238}_{92}$U is around $4.5\cdot 10^9$ years. There is $1$ gram of pure $^{238}_{92}$U atoms, how long does it take for it to have at least $0.75$ grams of them decay to some other substance.\footnote{The decay of radioactive substance is a random process: A standard nuclear reactors uses $20$ tonnes of Uranium a year, which is a scale of $10^{28}$ many Uranium atoms. Therefore, even though the half-life of Uranium seems to be very long, it is still highly likely to see some decay events happening at every single second given the large amount of atoms have in the reactor.}


\end{document}
